\section{Protected endpoints with at least three symbols}
\label{sec:nonbinary}
The logarithm in the corruption rate comes from terminal states with few
previous visits.  It is not an artifact of a compressed representation.
We first bound the perturbation of a direct full-history estimator, and
then construct an erasure channel whose \emph{entire output word} has only
a controlled amount of parameter information.
Throughout this section $K\geq3$, $M\geq64$, $0\leq\beta\leq1$, and
$k\in\{0,\ldots,M\}$.  Put $b_0=\max\{\beta,M^{-1}\}$,
$s=M^{-1/2}$, and
\[
 F_3=\begin{cases}
 0,&k=0,\\
 \displaystyle\frac{k}{M}\left[1+\log\left(1+\frac1{b_0k}\right)\right],&k\geq1.
 \end{cases}
\]

\subsection{Pathwise stability of corrupted counts}
Write $C_{abc}$ for the number of clean triples $(a,b,c)$ ending at
positions $3,\ldots,M+2$, and $C'_{abc}$ for the corresponding corrupted
counts.  A replacement affects at most three triple locations.  At each
location, the old and new indicator vectors differ in $\ell_1$ norm by at
most two.  Consequently, even with adjacent or simultaneous edits,
\begin{equation}\label{eq:nb-histogram}
 \|C'-C\|_1\leq6k.
\end{equation}
The terminal context $S=(X_{M+1},X_{M+2})$ is preserved.  Set
$J=\sum_cC_{Sc}$ and let
\[
 q_c=\frac{C_{Sc}+1}{J+K},\qquad
 q'_c=\frac{C'_{Sc}+1}{\sum_dC'_{Sd}+K}.
\]
For nonnegative vectors $u,v$ of positive masses $A,B$, insertion of
$v/A$ and the inequality $|A-B|\leq\|u-v\|_1$ give
\[
 \operatorname{TV}(u/A,v/B)
 \leq\tfrac12\bigl(\|u-v\|_1/A+|A-B|/A\bigr)
 \leq\|u-v\|_1/A.
\]
Apply this to the two smoothed rows, using the \emph{clean} mass in the
denominator, to obtain the deterministic bound
\begin{equation}\label{eq:nb-row-stability}
 \operatorname{TV}(q',q)\leq\min\{1,6k/(J+K)\}.
\end{equation}
This argument allows arbitrary changes to row membership.  It does not
condition row successors on the selected endpoint or treat them as IID.

We use the following inherited analytic interfaces.  They are stated here
to delimit the dependence on the predecessor dossier; the new corruption
argument does not re-prove its clean statistical theorem or its stationary
reverse-return analysis.  The R011 ORIGINAL rule is exactly the clean
add-one vector $q$ above.  Its uniform consequence applies for every
$K\geq2$, integer $M\geq1$, $\beta\in[0,1]$, every qualifying source,
and every initial pair law (including stationary initialization):
\begin{equation}\label{eq:nb-clean-interface}
 \mathbb E\operatorname{TV}(q,\theta_S)
 \leq2^{13}K^{5/2}M^{-1/2}.
\end{equation}
R011 gives a finer parameter-dependent bound, but its full collection
of scales is not needed here: the displayed uniform consequence is the
only clean upper interface used in this manuscript.  Its $K=2$ domain
also covers the endpoint result in Section~\ref{sec:endpoint}.

For a pair-chain path $Y_0,\ldots,Y_m$, define its actual-terminal
occupation count by
$J_m=\sum_{i=0}^{m-1}\mathbf1\{Y_i=Y_m\}$.
The IU001 stationary interface, for $\beta>0$, $m\geq1$, $A>0$, and
integer $\kappa\geq1$, is
\begin{equation}\label{eq:nb-IU-interface}
 \mathbb E_{\Pi}\min\{1,A/(J_m+\kappa)\}
 \leq\frac{dA}{m}\left[17+\log\left(1+\frac{480}{\beta A}\right)\right],
 \qquad d=|\operatorname{supp}\Pi|\leq K^2.
\end{equation}
This is a terminal-query bound under the two-step all-state certificate,
not an independent stationary-query identity.  The AU001 universal
interface, valid for every initial law, every $\beta\in[0,1]$, and the
same domains of $m,A,\kappa$, is
\begin{equation}\label{eq:nb-AU-interface}
 \mathbb E_{\nu}\min\{1,A/(J_m+\kappa)\}
 \leq\Psi\left(\frac{K^4A}{m+1}\right),\qquad
 \Psi(x)=\begin{cases}x(1+\log(1/x)),&0<x<1,\\0,&x=0,\\1,&x\geq1.
 \end{cases}
\end{equation}
For clarity, the short universal occupation argument can be reproduced.
For a chain on $d$ states and integer $w\geq1$, put
$b_m(y)=\mathbb P_y(J_m<w)$ and $B_m=\max_y b_m(y)$.
The count over a whole path dominates the count over an ending suffix,
so the Markov property gives $B_{m+1}\leq B_m$.  On each path,
$\sum_{t=0}^m\mathbf1\{J_t<w\}\leq dw$: each state can be encountered
with fewer than $w$ previous visits at most $w$ times.  Thus
\[
 (m+1)B_m\leq\sum_{t=0}^m B_t
 \leq\sum_y\sum_{t=0}^m b_t(y)\leq d^2w.
\]
Layer cake proves \eqref{eq:nb-AU-interface}.  Indeed, for $0<t<1$ the
strict threshold $J_m<A/t-\kappa$, when nonempty, is represented by
$w=\lceil A/t-\kappa\rceil\leq A/t$; integrate
$\min\{1,K^4A/((m+1)t)\}$.  This also explains why the smoothing
integer and the strict inequality matter.

\subsection{An exact suffix mixture for arbitrary initialization}
Assume first $\beta>0$ and $k\geq1$.  Let
\[
 u=\lfloor M/2\rfloor,\quad m=\lceil M/2\rceil,\quad
 r=\lfloor M/4\rfloor,\quad \eta=(1-\beta)^r.
\]
For $0<\beta<1$, the matrix
$R=(P^2-\beta\mathbf1\Pi)/(1-\beta)$ is stochastic on \emph{all}
pair states, and $\Pi R=\Pi$.  Expanding the powers of
$P^2=\beta\mathbf1\Pi+(1-\beta)R$ gives
\begin{equation}\label{eq:nb-suffix-mixture}
 \mu P^u=(1-\eta)\Pi+\eta\nu,
 \qquad \nu=\mu R^rP^{u-2r}.
\end{equation}
Extending each initial distribution with the ordinary transition matrix
$P$ gives the same exact mixture for the \emph{entire suffix path}
$Y_u,\ldots,Y_M$.  In particular, the residual suffix uses $P$, not $R$.
For $\beta=1$, $P^2=\mathbf1\Pi$ and the residual term is absent.
The full terminal count dominates the suffix terminal count pathwise.
Apply \eqref{eq:nb-IU-interface} and \eqref{eq:nb-AU-interface} with
$A=6k$, $\kappa=K$, and $m\geq M/2$:
\begin{align}\label{eq:nb-suffix-bound}
 \mathbb E\min\{1,6k/(J+K)\}
 &\leq12K^2\frac{k}{M}
 \left[17+\log\left(1+\frac{80}{\beta k}\right)\right]
 +\eta\Psi(12K^4k/M).
\end{align}
No synthetic regeneration event is supplied to the estimator; this
mixture is only an analysis of its existing path distribution.

For $\beta\geq1/M$, write $z=M/k\geq1$, $t=M\beta\geq1$,
$D=12K^4$, and $G=1+\log(1+z/t)$, so $F_3=G/z$.
We have $\eta\leq1$ and $\eta\leq e\exp(-t/4)$.
If $z<D$, then $\eta\Psi(D/z)\leq1\leq DG/z$.
If $z\geq D$, use
$\log z\leq\log(1+z/t)+\log t$ to obtain
\[
 \eta\Psi(D/z)\leq\frac Dz\eta(1+\log z)
 \leq\frac Dz\{G+e\exp(-t/4)\log t\}
 \leq\frac{5DG}{z}.
\]
The last elementary bound uses $e\exp(-t/4)\log t\leq4$ for $t\geq1$.
Also $\log(1+80x)\leq\log80+\log(1+x)$.
For $0\leq\beta<1/M$, apply \eqref{eq:nb-AU-interface} directly to
the full path.  Splitting at $z=6K^4$ proves
\[
 \Psi\left(\frac{6K^4k}{M+1}\right)
 \leq\Psi(6K^4/z)
 \leq\frac{6K^4}{z}[1+\log(1+z)]=6K^4F_3.
\]
Combining these estimates with \eqref{eq:nb-row-stability} and
\eqref{eq:nb-clean-interface} gives
\begin{equation}\label{eq:nb-final-upper}
 \mathbb E\operatorname{TV}(q',\theta_S)
 \leq\min\left\{1,2^{13}K^{5/2}s+
 [12K^2(17+\log80)+60K^4]F_3\right\}.
\end{equation}
At $k=0$ the perturbation vanishes and only the clean interface is used.
The rule $q'$ itself uses neither $k$ nor $\beta$.

\subsection{Exact erasure of arrangements}
\begin{lemma}[Exchangeable count erasure]\label{lem:nb-erasure}
Let $V_1,\ldots,V_\ell$ be IID Bernoulli$(p)$ bits, set
$n=\sum_iV_i$, fix an integer $b\geq1$, and put
$r=b\lfloor n/b\rfloor$.  Choose uniformly $n-r$ of the one positions
and replace them by zero.  The output $Z$, conditional on $r$, is uniform
among weight-$r$ words, independently of $p$.  At most $b-1$ bits are
changed.  If two further IID Bernoulli$(p)$ bits $U$ are untouched, the
joint experiment $(Z,U)$ factors through $(r,U)$ followed by a
$p$-independent channel, with at most
$4(\lfloor\ell/b\rfloor+1)$ messages.
\end{lemma}
\begin{proof}
Conditional on $n$, the original word is uniform among its
$\binom\ell n$ arrangements.  A fixed output of weight $r$ is contained
in $\binom{\ell-r}{n-r}$ original one sets, and is chosen from each one
set with probability $1/\binom nr$.  Hence
\[
 \mathbb P(Z=z\mid n)
 =\frac{\binom{\ell-r}{n-r}}{\binom\ell n\binom nr}
 =\frac1{\binom\ell r}.
\]
Averaging within a cell of the rounded count leaves this identity
unchanged.  Conditional on $p$, the two untouched bits are independent
of the first $\ell$ bits and the erasure randomization.  Thus the same
uniform-arrangement channel works conditional on $(r,U)$.
The identity includes $\ell=0$ and the incomplete last count cell.
\end{proof}

\begin{lemma}[Finite-message Bayes loss]\label{nb:finite-message}
If a real parameter is uniform on an interval of length $h>0$ and an
arbitrary stochastic channel emits at most $D$ messages, every
independently randomized real-valued decoder has expected absolute loss
at least $h/(4D)$.
\end{lemma}
\begin{proof}
Fix its independent decoding seed, obtaining at most $D$ possible output
points.  Their radius-$t$ neighborhoods have prior mass at most $2Dt/h$.
Whatever stochastic encoding is used, the event of error at most $t$
requires that the parameter lie in that union.  Integrate
$1-2Dt/h$ over $0\leq t\leq h/(2D)$ and then average the seed.
\end{proof}
In Lemma~\ref{lem:nb-erasure}, a full-word decoder can be composed with
the parameter-independent uniform-arrangement channel, so it is a
randomized decoder of the finite message.  This is not an assumption
that the observation is only a histogram.  The protected pair is part
of the message, and the prior is not asserted to remain uniform after
conditioning on that pair.  Neither the actual edit mask nor the actual
attack seed is disclosed.

\subsection{Two IID floors common to all mixing classes}
\begin{lemma}\label{nb:iid-floors}
For every $K\geq2$, $M\geq64$, and $\beta\in[0,1]$, the stationary
protected-endpoint minimax risk is at least $s/64$.  For $k\geq1$ it is
also at least $k/(96M)$.  Each lower bound is realized by a stationary
IID binary family and a common corruption kernel across that family.
\end{lemma}
\begin{proof}
Choose two named symbols $C_0,C_1$ and set every source row, including
unused contexts, to $(1-p)\delta_{C_0}+p\delta_{C_1}$.  Its two-step
pair matrix has every row equal to the stationary pair law, so it
satisfies every certificate under consideration.
For the statistical floor use identity corruption and
$p_\pm=1/2\pm1/(16\sqrt N)$, $N=M+2$.
The one-sample divergence is at most $32/(256N)$, hence the product
KL divergence is at most $1/8$ and Pinsker's inequality bounds product
TV by $1/4$.  The target laws are separated in TV by $1/(8\sqrt N)$.
Nearest-target testing and the triangle inequality give average loss at
least one quarter of this separation times $1-1/4$, namely
$3/(128\sqrt N)\geq s/64$.

For the corruption floor use one prior uniform on $[1/4,3/4]$ and
$b=k+1$.  On all-binary input words apply Lemma~\ref{lem:nb-erasure}
to the first $M$ bits and leave the last two untouched; on every other
input use identity.  This defines a common map on the entire alphabet
with at most $k$ edits.  The full output has at most
$D=4(\lfloor M/b\rfloor+1)$ informative messages.
Since TV loss dominates the absolute error in the $C_1$ coordinate,
Lemma~\ref{nb:finite-message} gives
\[
 R_{\mathrm{stat}}^{\mathrm{prot}}
 \geq\frac1{32(\lfloor M/b\rfloor+1)}
 \geq\frac b{32(M+b)}\geq\frac{k}{96M},
\]
where $b\leq M+1\leq2M$.  The statistical and corruption witnesses
need not coincide: their lower bounds are combined by a maximum.
\end{proof}

\subsection{Stationary rare suffixes and the full logarithm}
Suppose $b_0\leq1/16$.  Put $a=1/(2M)$ and $e=b_0-a$.
For $p\in[1/4,3/4]$, $p_0=1-p$, $p_1=p$, use symbols $W,C_0,C_1$
and define a first-order chain by
\begin{align*}
 T(W,W)&=1-a,& T(W,C_i)&=ap_i,\\
 T(C_i,W)&=e,&T(C_i,C_j)&=(1-e)p_j.
\end{align*}
Its stationary symbol law is
\[
 \pi(W)=e/b_0,\qquad\pi(C_i)=ap_i/b_0.
\]
Give unused symbols zero stationary mass and set every unused starting
row to $\pi$.  Direct substitution verifies $\pi T=\pi$.
The coordinate inequalities $T(y,v)\geq b_0\pi(v)$ hold for every
starting symbol $y$, including unused ones: the $W$-to-active and
active-to-$W$ ratios equal $b_0$, while the other inequalities reduce
to $b_0\leq1$.  Set $\theta_{xy}(v)=T(y,v)$ on \emph{every} pair and
$\Pi(u,v)=\pi(u)T(u,v)$.  Then $\Pi$ is stationary and
\begin{equation}\label{eq:nb-allstate-witness}
 P^2((x,y),(u,v))=T(y,u)T(u,v)
 \geq b_0\Pi(u,v)\geq\beta\Pi(u,v)
\end{equation}
on all pair states, not merely the stationary support.
Initialize at this stationary pair law, which is not $\pi\otimes\pi$.

Let $L=\lfloor1/(4b_0)\rfloor$, so $4\leq L\leq M/4$.
For $2\leq j\leq L$, let $E_j$ be the event that the complete word is
$W^{M+2-j}$ followed by $j$ active letters.  These disjoint events have
probabilities
\[
 w_j=\Pi(W,W)(1-a)^{M-j}a(1-e)^{j-1}.
\]
Here $\Pi(W,W)=(1-a/b_0)(1-a)\geq1/4$ is independent of $p$.
Bernoulli's inequality gives
$(1-a)^{M-j}\geq1/2$ and
$(1-e)^{j-1}\geq1-b_0L\geq3/4$.  Thus
\begin{equation}\label{eq:nb-event-mass}
 w_j\geq3/(64M)\geq1/(32M).
\end{equation}
Conditional on $E_j$, the $j$ active bits are IID Bernoulli$(p)$.
The same uniform prior on $p$ is retained after this event, before the
endpoint labels are observed.

Define a single kernel on all words as follows.  If its input is
$W^{N-j}$ followed by $j$ active symbols with $2\leq j\leq M$, round
the number of $C_1$ symbols among the first $j-2$ active symbols to a
multiple of $b=k+1$, using the uniform deletion rule.  Preserve all
markers and the final pair.  Use identity on every other word.
This kernel depends on neither $p$ nor $\beta$, preserves the visible
suffix length $j$, and makes at most $k$ replacements in allowed
positions on every input.  It is the same kernel for every suffix and
every parameter value.

At all active terminal pairs the actual future law is
\begin{equation}\label{eq:nb-true-future}
 e\delta_W+(1-e)[(1-p)\delta_{C_0}+p\delta_{C_1}].
\end{equation}
The return mass $e$ must not be omitted merely because no return
occurred in the selected past.  For every prediction $q$,
$\operatorname{TV}(q,\theta_S)\geq|q(C_1)-(1-e)p|$.
Apply Lemmas~\ref{lem:nb-erasure} and \ref{nb:finite-message} to
$q(C_1)/(1-e)$, with interval length $1/2$.  Since $1-e\geq1/2$,
the conditional Bayes loss is at least
\[
 \frac1{64[1+(j-2)/b]}.
\]
The decoder may be granted the value of $j$, which is already visible.
The output arrangement adds no information beyond the finite message.
Sum the losses over the disjoint events, without hiding observations
outside them:
\begin{align}\label{eq:nb-harmonic-lower}
 R_{\mathrm{stat}}^{\mathrm{prot}}
 &\geq\frac1{2048M}\sum_{\ell=0}^{L-2}\frac1{1+\ell/b}
 \geq\frac b{2048M}\log\left(1+\frac{L-1}{b}\right).
\end{align}
The last inequality integrates a decreasing function over $[0,L-1]$.
Since $L-1\geq1/(8b_0)$, the function
$b\log(1+c/b)$ is increasing in $b>0$, and
$\log(1+x/8)\geq\log(1+x)/8$, it follows that
\begin{equation}\label{eq:nb-log-lower}
 R_{\mathrm{stat}}^{\mathrm{prot}}
 \geq\frac{k}{16384M}\log\left(1+\frac1{b_0k}\right).
\end{equation}
The same prior and same kernel were used across all $j$; choosing a
different source separately at each suffix would not justify this sum.

For $b_0\leq1/16$, maximize \eqref{eq:nb-log-lower} and the IID
$k/(96M)$ floor to obtain $F_3/32768$.  For $b_0>1/16$,
$\log(1+1/(b_0k))\leq\log17<3$, so $F_3\leq4k/M$ and the IID floor
alone is stronger than $F_3/32768$.  Maximizing with the statistical
floor gives
\begin{equation}\label{eq:nb-final-lower}
 R_{\mathrm{stat}}^{\mathrm{prot}}
 \geq2^{-15}\min\{1,\max(s,F_3)\}.
\end{equation}
For $k=0$, use only the statistical floor.  Together with
\eqref{eq:nb-final-upper} and inclusion of stationary initializations
among arbitrary ones, this proves the claimed nonbinary profile.

\subsection{The zero-certificate mechanism without a mixing cutoff}
The preceding construction already covers $\beta=0$, but a simpler
version illuminates the logarithm and records the original stationary
strengthening.  This paragraph is valid for every $M\geq2$, not only
$M\geq64$.
For arbitrary initialization, start at $WW$, set $\alpha=1/M$, and use
\[
 T(W,W)=1-\alpha,\quad T(W,C_i)=\alpha p_i,
 \quad T(C_i,C_j)=p_j.
\]
Unused rows can equal $\delta_W$.  A qualifying stationary law is
supported on the active IID class, so the source belongs to
$\mathcal D_2(0)$ despite its nonstationary start.  For
$2\leq j\leq M$, the event $E_j$ above has mass
$\alpha(1-\alpha)^{M-j}\geq1/(4M)$, independent of $p$.
The future active law is exactly Bernoulli$(p)$.
The same marker erasure gives conditional risk at least
$1/[32(1+(j-2)/(k+1))]$.  Hence, for $b=k+1$,
\[
 R_{\mathrm{all}}^{\mathrm{prot}}
 \geq\frac b{128M}\log\left(1+\frac{M-1}{b}\right)
 \geq\frac1{512}\Psi(k/M),\qquad k\geq1.
\]
For the last inequality, use $M-1\geq M/2$, monotonicity in $b$, and
$\log(1+t/2)\geq(1+\log t)/4$ for $t=M/k\geq1$.
The latter follows by checking $t=1$ and differentiating.

For a stationary version set $a=e=1/M$ in the recurrent marker chain.
Now $\pi(W)=1/2$, $\pi(C_i)=p_i/2$, and the initialization is
$\Pi(u,v)=\pi(u)T(u,v)$.  For every $2\leq j\leq M$,
\[
 \mathbb P_p(E_j)
 =\frac{(1-1/M)^M}{2M}\geq\frac1{8M}.
\]
The equality includes the initial $WW$ factor; it is independent of
both $p$ and $j$.  The lower bound follows from
$(1-1/M)^M\geq1/4$ for $M\geq2$.
The true future law retains its return mass $e$, reducing the
conditional lower by at most a factor two.  The same summation yields
$\Psi(k/M)/2048$.  Thus transience is not needed for the logarithmic
obstruction.  No concentration cutoff such as $j=O(k^2)$ is needed:
exact arrangement erasure works for every suffix length.

\section{When the original terminal context is not protected}
\label{sec:endpoint}
We now change exactly one part of the experiment.  The adversary may
replace any of the $N=M+2$ positions, including the final two, subject
to the same strict budget $k\leq M$.  The prediction target remains
$\theta_{X_{N-1},X_N}$ at the \emph{original hidden} terminal pair.
It is not the row indexed by the corrupted endpoint.  Write
$R_{\mathrm{stat}}^{\mathrm{end}}$ and
$R_{\mathrm{all}}^{\mathrm{end}}$ for these risks and put
\[
 \delta=1-\beta,\qquad
 H_{\mathrm{end}}=\min\{1,\max(s,k/M,\delta\mathbf1\{k\geq1\})\}.
\]

\begin{proposition}\label{prop:endpoint}
For $K\geq2$, $M\geq64$, $0\leq\beta\leq1$, and integer
$0\leq k\leq M$,
\[
 H_{\mathrm{end}}/96\leq R_{\mathrm{stat}}^{\mathrm{end}}
 \leq R_{\mathrm{all}}^{\mathrm{end}}
 \leq C_K^{\mathrm{end}}H_{\mathrm{end}},
 \qquad
 C_K^{\mathrm{end}}=\max\{2^{13}K^{5/2},\sqrt K/2+3\}.
\]
For $k=0$ the clean selected-row add-one rule suffices.  For $k\geq1$
the pooled empirical distribution of corrupted symbols at positions
$3,\ldots,N$ suffices.  Neither branch requires $\beta$.
\end{proposition}

\subsection{What the all-state certificate says about each row}
Let $\pi(c)=\sum_d\Pi(c,d)$.  Summing
$P^2(u,(c,d))\geq\beta\Pi(c,d)$ over $d$ gives
\begin{equation}\label{eq:end-row-domination}
 \theta_u(c)\geq\beta\pi(c)
 \quad\hbox{for every pair $u$ and symbol $c$}.
\end{equation}
Indeed the first coordinate after two pair transitions is precisely the
next generated symbol from $u$.  The conclusion includes contexts
outside stationary support.  For $\beta<1$, write
$\theta_u=\beta\pi+\delta r_u$, where $r_u$ is a probability vector.
For $\beta=1$, coordinatewise domination and unit mass imply
$\theta_u=\pi$.  In every case,
\begin{equation}\label{eq:end-row-tv}
 \operatorname{TV}(\theta_u,\pi)\leq\delta.
\end{equation}
At $\beta=1$ all generated symbols from position 3 onward are IID
$\pi$, even for an arbitrary initial pair; no assertion about that
initial pair being IID is needed.

\subsection{A pooled martingale upper bound}
For $k\geq1$, let
$\bar p=M^{-1}\sum_{t=3}^N e_{X_t}$ be the clean generated-symbol
empirical distribution.  With the natural filtration, define
$m_t=\theta_{X_{t-2},X_{t-1}}$ and $D_t=e_{X_t}-m_t$.
These are vector martingale differences.  Orthogonality at distinct
times gives
\[
 \mathbb E\left\|\frac1M\sum_{t=3}^ND_t\right\|_2^2
 =\frac1{M^2}\sum_{t=3}^N\mathbb E[1-\|m_t\|_2^2]
 \leq\frac1M.
\]
Cauchy--Schwarz and Jensen imply that the expected half-$\ell_1$ norm
of this martingale average is at most $\sqrt K/(2\sqrt M)$.
By \eqref{eq:end-row-tv} and convexity, the average $M^{-1}\sum_tm_t$
is within TV $\delta$ of $\pi$ on every path, as is the true terminal
next law.  A pathwise triangle inequality therefore yields
\[
 \mathbb E\operatorname{TV}(\bar p,\theta_{X_{N-1},X_N})
 \leq\frac{\sqrt K}{2\sqrt M}+2\delta.
\]
Let $\bar p'$ pool the corrupted symbols at positions $3,\ldots,N$.
Each changed symbol alters this empirical distribution by at most
$1/M$ in TV, so $\operatorname{TV}(\bar p',\bar p)\leq k/M$ pathwise.
Edits to the first two positions contribute nothing to this bound.
Thus
\begin{equation}\label{eq:end-upper}
 \mathbb E\operatorname{TV}(\bar p',\theta_{X_{N-1},X_N})
 \leq\min\{1,\sqrt K s/2+2\delta+k/M\}.
\end{equation}
This proof does not reconstruct the endpoint, appeal to mixing-time
concentration, or substitute an independent query for the actual one.
For $k=0$, use \eqref{eq:nb-clean-interface} instead: the pooled bias
$2\delta$ is only an upper-bound artifact in the clean problem.
These two branches prove the upper in Proposition~\ref{prop:endpoint}.

\subsection{A one-symbol erasure lower bound even for a known source}
Assume $k\geq1$.  Set $\alpha=\min\{\delta/2,1/4\}$ and consider the
symmetric first-order binary chain
\[
 T(0,0)=T(1,1)=\tfrac12+\alpha,
 \qquad T(0,1)=T(1,0)=\tfrac12-\alpha.
\]
Its stationary symbol law is uniform on $\{0,1\}$.  For $K>2$, set
unused starting rows to that same binary uniform law and all unused
stationary masses to zero.  Extend to every second-order row by
$\theta_{xy}(c)=T(y,c)$ and initialize at
$\Pi(u,v)=\pi(u)T(u,v)$.  Every starting symbol satisfies
$T(y,u)\geq(1-2\alpha)\pi(u)$, including unused symbols, so
\[
 P^2((x,y),(u,v))=T(y,u)T(u,v)
 \geq(1-2\alpha)\Pi(u,v)\geq\beta\Pi(u,v).
\]
The witness therefore obeys the required all-state certificate.

Use the deterministic common attack that sets only the last symbol to
$0$ and leaves the rest unchanged.  It makes at most one replacement
on every word and requires no source parameters.  Grant the estimator
the entire source law and its stationary initialization.  The output
then reveals the true prefix through $X_{N-1}$ but not $X_N$.
Conditionally on that prefix, the probabilities of the two possible
last bits are $1/2+\alpha$ and $1/2-\alpha$ in some order.  The two
possible true next-law probabilities of symbol $1$ are
$1/2-\alpha$ and $1/2+\alpha$.
For any real prediction, the least posterior absolute error is the
smaller last-bit probability times the gap between these target values:
\[
 (\tfrac12-\alpha)\,2\alpha=\alpha-2\alpha^2.
\]
To check this weighted-median formula, the objective is linear below,
between, and above the two target values, and its minimum occurs at
the more likely endpoint.  Randomization cannot lower a linear average
of deterministic risks.  TV loss dominates the absolute error in the
symbol-$1$ coordinate, including predictions assigning mass to unused
symbols.  The bound holds after every prefix and hence unconditionally.

If $\delta\leq1/2$, the loss is
$\delta(1-\delta)/2\geq\delta/4$; if $\delta\geq1/2$, it is
$1/8\geq\delta/8$.  Therefore
\begin{equation}\label{eq:end-one-edit}
 R_{\mathrm{stat}}^{\mathrm{end}}\geq\delta/8,
 \qquad k\geq1.
\end{equation}
This obstruction is uncertainty about a latent final symbol in a
\emph{known} source, rather than uncertainty in estimating its rows.
The model does not reveal a replacement indicator: for this particular
attack that indicator would itself reveal the erased binary symbol.
No conclusion about all possible visible-mask models follows.

\subsection{Assembly and scope of the endpoint contrast}
Both IID floors in Lemma~\ref{nb:iid-floors} remain valid.  Their
identity and first-$M$-symbol erasure kernels are allowed in the enlarged
corruption class, even though they do not use endpoint edits.
For $k\geq1$, maximize $s/64$, $k/(96M)$, and
\eqref{eq:end-one-edit}; this gives
$\max(s,k/M,\delta)/96$, and hence the required capped lower.
For $k=0$, only the statistical floor is needed.
All lower witnesses are stationary, and the inclusion of stationary
initializations among arbitrary ones completes the proof.

In particular, at any fixed $\beta<1$, one allowed endpoint edit
prevents consistency for the original-terminal target.  At $\beta=1$
all source rows coincide and the endpoint ambiguity disappears.
For positive budgets, statistical-order risk is possible only when
$1-\beta=O(M^{-1/2})$ and $k=O(\sqrt M)$, and these two conditions
are also sufficient up to fixed-$K$ constants by the proved bounds.
This comparison concerns full-history statistical prediction.  It does
not assert an online adversarial equivalence, a visible-mask result,
a neural-depth requirement, or a training guarantee.
